% ECONOMICS_PROBLEM: {"id": "O31-378", "title": "Research productivity versus commercialization", "jel_code": "O31", "source_id": "OP-0261"}
% !TeX program = xelatex
\documentclass[11pt,a4paper]{article}
\usepackage[margin=23mm]{geometry}
\usepackage{fontspec}
\setmainfont{Latin Modern Roman}
\usepackage{amsmath,amssymb,mathtools}
\usepackage{xcolor,enumitem,booktabs,longtable,array,xurl,hyperref}
\definecolor{muted}{HTML}{53636C}
\hypersetup{colorlinks=true,urlcolor=blue,linkcolor=blue}
\setlength{\parindent}{0pt}
\setlength{\parskip}{5pt}
\newcommand{\R}{\mathbb R}
\newcommand{\E}{\mathbb E}
\newcommand{\N}{\mathbb N}
\newcommand{\ind}{\mathbf 1}
\newcommand{\Var}{\operatorname{Var}}
\newcommand{\Cov}{\operatorname{Cov}}
\newcommand{\supp}{\operatorname{supp}}
\DeclareMathOperator*{\argmin}{arg\,min}
\DeclareMathOperator*{\argmax}{arg\,max}
\newcommand{\entrymeta}[3]{{\sffamily\small\color{muted}\textbf{\texttt{#1}}\quad|\quad #2\par #3\par}}
\newcommand{\Problemref}[1]{\texttt{#1}}
\newcommand{\JELref}[2]{\texttt{#2} (JEL \texttt{#1})}
\begin{document}
\section*{Research productivity versus commercialization}
\textbf{Problem ID:} \texttt{O31-378}\quad \textbf{JEL:} \texttt{O31}\par
% BEGIN ECONOMICS STATEMENT
\label{problem:OP-0261}
{\sffamily\small\color{muted}Original subject group: Growth and productivity\par}
\label{definitions:problem:30007542}
\textbf{Audit classification:} Published research agenda. The abstract invites renewed ideas/growth research; separating three latent blocks is an editorial specification. Verification concerns the cited version, not first priority or proof that this exact editorial specification remains unsolved on October 8, 2026.
\subsection*{Definitions and precise research target}
\paragraph{Editorial mathematical specification.} Fix industries \(j\) and years \(t\). Let \(R_{jt}\) measure research input, \(I_{jt}\) latent economically useful new ideas, \(P_{jt}\) observed patents adjusted by a predeclared quality measure, and \(A_{jt}\) measured productivity. Specify a research relation \(I_{jt}=F_\theta(R_{jt},A_{jt},v_{jt})\), a patent measurement relation \(P_{jt}=M_\psi(I_{jt},z_{jt})+\epsilon_{jt}\), and an adoption relation \(A_{j,t+h}=D_\gamma(I_{\le t},K_{jt},b_{jt},h)+\nu_{j,t+h}\). Here \(v\) denotes research opportunities, \(z\) patenting incentives, \(K\) complementary capital, and \(b\) commercialization barriers.
Identify changes in research productivity \(\partial F_\theta/\partial R\) separately from changes in patent measurement and adoption. The target is a decomposition of a declared productivity-growth slowdown into research-output and implementation components, with interactions allocated by an explicitly stated counterfactual convention. Use variation affecting research effort and adoption costs differently, explaining exclusion restrictions and spillovers across firms. Test whether competing parameterizations fit both patent-quality distributions and independently measured product or process improvements. Report identified sets when latent idea quality prevents point identification. Patent counts alone cannot establish that ideas become uniformly harder to find, while rising patents alone cannot establish increasing economically relevant research productivity.

Latent idea units require a normalization: for example fix the loading of a declared independently validated invention-quality proxy to one, and model a second proxy with independent errors conditional on quality. Otherwise scaling \(I\) and compensating \(F,M,D\) leaves observations unchanged, so marginal research productivity in unspecified idea units is unidentified. Research inputs must include real researcher effort and complementary resources, not only expenditure inflated by research wages. Specify adoption lags and the no-new-ideas counterfactual, holding initial knowledge and other productivity determinants fixed. Define the compatible set of triples \((\theta,\psi,\gamma)\) by the joint law of research, patents, independent innovation proxies, and output. A slowdown decomposition needs a fourth residual/other-output component unless the three structural blocks exhaust productivity changes.
\subsection*{Variants and assumption changes}
The following are \emph{editorial variants}, not additional published open conjectures.
\begin{enumerate}
\item \emph{Patent production only.} Treat quality-adjusted patents as the explicitly named research output and identify research-to-patent elasticities using valid effort variation. This narrower object can be estimated without claiming that patents measure all useful ideas.
\item \emph{Commercialization and diffusion.} Introduce separately observed product launches and adoption among noninventing firms, with research held fixed. Identify the lag from invention to productivity and bound the contribution of unpatented innovation.
\end{enumerate}
\subsection*{References and source audit}
\begin{itemize}
\item Teresa C. Fort; Nathan Goldschlag; Jack Liang; Peter K. Schott; Nikolas Zolas. \emph{Growth is Getting Harder to Find, Not Ideas}. 2025. Locator: Census CES25-21, April 2025, Abstract final paragraph. Primary record and abstract/summary checked; no fulltext claim is made. \url{https://www.census.gov/library/working-papers/2025/adrm/CES-WP-25-21.html}.
\end{itemize}
The abstract invites renewed ideas/growth research; separating three latent blocks is an editorial specification.
\begin{enumerate}\raggedright
\item\hypertarget{definitions:source:30007542:1}{} Teresa C. Fort; Nathan Goldschlag; Jack Liang; Peter K. Schott; Nikolas Zolas. 2025. \emph{Growth is Getting Harder to Find, Not Ideas}. Census CES25-21, April2025, Abstract final paragraph.
\url{https://www.census.gov/library/working-papers/2025/adrm/CES-WP-25-21.html}.
\end{enumerate}

\subsection*{Common conventions: Source mathematical conventions}

These conventions apply to each entry unless explicitly replaced.

\textbf{Models and identification.} A primitive model \(m\) specifies all probability laws, feasible choices, information, equilibrium and measurement rules used by its target. The admissible class \(\mathcal M\) is fixed independently of outcomes. If \(P_O(m)\) is its observable probability law and \(\tau(m)\) the target, its identified set at observable law \(P\) is
\[
 \mathcal I_\tau(P)=\{\tau(m):m\in\mathcal M,\ P_O(m)=P\}.
\]
Point identification means that this set is a singleton. An empty set signals incompatibility of data and assumptions. Coordinate bounds need not describe the joint set. Bounds are sharp only if every retained target is attainable or arbitrarily approachable by compatible models. Finite bounds require explicit restrictions on unobserved quantities; unrestricted confounding, hidden wealth, extrapolation or arbitrary equilibrium selection can yield uninformative sets.

\textbf{Causal interventions.} A potential outcome \(Y_i(p)\) is the outcome under a complete policy regime \(p\), including timing, financing, information and market spillovers. The target population and units are fixed before treatment. With interference, use the complete assignment vector or a justified exposure mapping; individual no-interference notation alone cannot identify an economy-wide intervention. A difference of mean potential outcomes does not require the cross-world joint distribution. Conditioning on post-treatment migration, survival, enrollment or employment changes the estimand and needs separate justification.

\textbf{Statistical statements.} An estimator, confidence set, forecast or test specifies a sampling law, model class and loss/coverage criterion. Uniform \((1-\alpha)\)-coverage means \(\inf_{m\in\mathcal M}\Pr_m\{\tau(m)\in C_n\}\geq1-\alpha\), or an explicitly asymptotic version. The whole parameter space trivially has coverage, so informativeness requires a stated width, risk or power objective. A forecasting improvement is relative to a named benchmark, information set and evaluation population.

\textbf{Equilibrium and optimization.} An equilibrium comprises feasible plans, optimality, beliefs where needed, budget constraints and all market/resource clearing. The primitives or a stated benchmark define each correspondence \(\mathcal E(p;m)\). Nonemptiness is assumed or proved, never inferred just from compact policy sets. A selected-equilibrium target uses a declared selection rule; a robust target quantifies over all permitted equilibria. Use suprema or infima unless attainment is established. An \(\varepsilon\)-optimal policy is feasible and within \(\varepsilon\) of the finite optimum value. Report an empty feasible set separately.

\textbf{Welfare and accounting.} Normative utility, interpersonal normalization, social weights, numeraire, discounting and the use of public receipts are primitives. Behavioral utility need not coincide with the welfare criterion. Household welfare includes ownership income and rents once; adding profits again to compensating variation generally double-counts them. A monetary equivalent is defined by an explicit transfer and price convention, with monotonicity and range conditions. Average effects, mechanism contributions, transfers and welfare losses are different targets. Infinite sums require convergence and dynamic feasible plans require terminal or no-Ponzi conditions.

\textbf{Source versions.} A working-paper date, revision date and journal publication year are distinguished. A DOI for a journal version is not automatically the DOI of an earlier manuscript. Locators refer to the stated version and printed pages where specified. A metadata-only check does not verify a claim about a full-text conclusion. Every entry's source subsection narrows the provenance claim accordingly.
% END ECONOMICS STATEMENT
\end{document}
